% ============================================================
% CAPITOLO 20 - I RITARDI
% ============================================================
\chapter{I ritardi}
\label{chap:ritardi}

Il \textbf{ritardo} è la più importante tra le figure di dissonanza melodica: una nota estranea che crea una tensione momentanea prima di risolvere. A differenza delle note estranee del capitolo precedente, il ritardo è sempre preparato.

\section{Definizione e meccanismo}
\label{sec:definizione_ritardo}

Il ritardo è una dissonanza sul tempo forte di una misura (o su una parte forte di un tempo) che ritarda l'arrivo di una nota reale dell'accordo. Si svolge in tre fasi, da rispettare con rigore.

\begin{definizione}[Le tre fasi del ritardo]
	\label{def:fasi_ritardo}
	Un ritardo si articola sempre in tre momenti, gli stessi del ciclo generale della dissonanza (Regola \ref{reg:trattamento_dissonanza}):
	\begin{enumerate}
		\item \textbf{Preparazione:} la nota che diventerà dissonante è introdotta come consonanza nell'accordo precedente, preferibilmente nella stessa voce e alla stessa altezza.
		\item \textbf{Percussione:} l'accordo cambia, ma la nota preparata è tenuta (legata) e diventa dissonante nel nuovo contesto. Avviene su un tempo forte.
		\item \textbf{Risoluzione:} la nota dissonante scende di grado congiunto su una nota reale dell'accordo.
	\end{enumerate}
\end{definizione}

\section{Il ritardo 4-3}
\label{sec:ritardo_4_3}

Il ritardo 4-3 è il più usato. Ritarda la terza dell'accordo: sulla dominante ritarda la sensibile (Si in Do maggiore), sulla tonica ritarda la terza (Mi in Do maggiore). La quarta è dissonante rispetto al basso, la terza è la sua risoluzione.

\begin{regola}[Cifratura del ritardo]
	\label{reg:cifratura_ritardo}
	Nel basso figurato il ritardo 4-3 si indica scrivendo in sequenza \texttt{4 3}. Il 4 rappresenta la dissonanza (la quarta sul basso), il 3 la sua risoluzione (la terza dell'accordo).
\end{regola}

Il caso più semplice si ha sulla dominante: la nota ritardata, il Do in Do maggiore, è la tonica, già presente nell'accordo che precede (I o IV). Viene tenuta sul basso Sol e scende alla sensibile, che durante il ritardo non compare in nessun'altra voce.

\begin{esempio}[Ritardo 4-3 sulla dominante]
	\label{ex:sospensione_43}
	Il Do del soprano, preparato dall'accordo di IV, è tenuto sul V (cifra 4) e scende al Si (cifra 3).
	\musica{es_cap11_04}
\end{esempio}

\section{Il ritardo nella cadenza perfetta}
\label{subsec:ritardo_cadenza}

Nella cadenza perfetta (V-I) il ritardo 4-3 amplifica la tensione della dominante. Il \sigla{I}{6}{4} cadenzale (Sezione \ref{sec:cadenziale}) è in realtà un \textbf{doppio ritardo} sull'accordo di dominante. Sul basso Sol, fermo, la sesta (Mi) scende alla quinta (Re) e la quarta (Do) scende alla terza (Si): le cifre \texttt{6-5} e \texttt{4-3} si risolvono insieme. Per questo le regole del \sigla{I}{6}{4} cadenzale coincidono con quelle del ritardo: le note in quarta e sesta sono preparate, cadono sul tempo forte e scendono per grado.

\section{Il ritardo nel modo minore}
\label{subsec:ritardo_minore}

Nel modo minore il ritardo 4-3 segue gli stessi principi, ma si adatta alla scala minore armonica: sulla dominante la nota ritardata (il La in La minore) scende al Sol$\sharp$, la sensibile alzata. L'esempio seguente mostra il ritardo nel soprano, in Do maggiore e in La minore.

\begin{esempio}[Ritardo 4-3 nel soprano]
	\label{ex:ritardo_4_3_s}
	\musica{es_cap10_4_3_s}
\end{esempio}

\section{Il ritardo nel soprano e nel contralto}
\label{sec:ritardo_4_3_voci}

Il ritardo 4-3 sulla dominante può essere affidato a una qualunque voce superiore. La nota ritardata è tenuta sul basso e scende di grado sulla sensibile, che durante il ritardo non compare in altre voci.

\begin{esempio}[Ritardo 4-3 nel contralto]
	\label{ex:ritardo_4_3_a}
	\musica{es_cap10_4_3_a}
\end{esempio}

\begin{errore}[Raddoppiare o anticipare la nota di risoluzione]
	\label{err:ritardo_risoluzione}
	Scrivere in un'altra voce, durante il ritardo, la nota su cui il ritardo risolve (per esempio il Si accanto al Do ritardato del 4-3) produce un urto di seconda minore o di nona, annulla l'effetto del ritardo e viene corretto come errore.
\end{errore}

\begin{regola}[Condotta del ritardo]
	\label{reg:condotta_ritardo_base}
	\begin{itemize}
		\item La nota ritardata è \textbf{la stessa} nella preparazione e nel ritardo (legatura di valore).
		\item Il ritardo cade sul tempo forte, la risoluzione sul debole, \textbf{sempre per grado discendente}.
		\item Durante il ritardo la \textbf{nota di risoluzione non compare in nessuna voce superiore}, per evitare la seconda (o nona) aspra.
		\item Si evitano le quinte e le ottave parallele.
	\end{itemize}
	I ritardi che riguardano la nona e la settima hanno ulteriori vincoli (Capitolo \ref{chap:ritardi_complessi}).
\end{regola}
